\section*{Bab \ref{ch:g12:projectile-motion} --- Jatuh Bebas dan Gerak Proyektil}

\begin{solution}{exo:g12:projectile-motion:1}
$h = \tfrac12 \times 9.81 \times 2.0^2 = \qty{19.6}{m}$;
$v = gt = \qty{19.6}{m/s} \approx \qty{71}{km/h}$.
\end{solution}

\begin{solution}{exo:g12:projectile-motion:2}
Tidak ada udara di Bulan: keduanya dalam \omterm{def:g12:projectile-motion:freefall}{jatuh bebas}, $m\vect a = m\vect g$
memberikan $\vect a = \vect g$ untuk massa berapa pun. Di Bumi, hambatan udara
pada bulunya sebanding dengan \omterm{def:g10:universal-gravitation:weight}{beratnya} yang mungil: bulunya melayang pada \omterm{def:g12:projectile-motion:terminal}{laju
terminal} sekitar \qty{1}{m/s} sementara palunya hampir tidak merasakan udaranya.
\end{solution}

\begin{solution}{exo:g12:projectile-motion:3}
$t = \sqrt{2h/g} = 0.45$, $1.4$, \qty{4.5}{s};
$v = gt = 4.4$, $14$, \qty{44}{m/s}. Waktunya tumbuh seperti $\sqrt h$:
$\times 100$ pada ketinggiannya hanya $\times 10$ pada waktunya.
\end{solution}

\begin{solution}{exo:g12:projectile-motion:4}
$h = v_0^2/2g = \qty{11.5}{m}$ pada $t = v_0/g = \qty{1.53}{s}$;
kembali pada \qty{3.06}{s}, pada \qty{15}{m/s} lagi (ketangkupannya).
\end{solution}

\begin{solution}{exo:g12:projectile-motion:5}
$t = \sqrt{2 \times 0.80/9.81} = \qty{0.40}{s}$;
$x = 2.5 \times 0.40 = \qty{1.0}{m}$. Tidak: waktu jatuhnya milik
gerak tegaknya saja (kebebasannya).
\end{solution}

\begin{solution}{exo:g12:projectile-motion:6}
$(\ang{15},\ang{75})$ dan $(\ang{30},\ang{60})$ berbagi \omterm{def:g12:projectile-motion:range}{jangkauannya};
\ang{45} menang. Yang lebih curam dari sepasang sudut berpelurus terbang lebih
lama: $t_f = 2v_0\sin\alpha/g$ tumbuh bersama $\alpha$.
\end{solution}

\begin{solution}{exo:g12:projectile-motion:7}
$4.905\,t^2 + 5.0\,t - 20 = 0$ memberikan
$t = (-5 + \sqrt{25 + 392.4})/9.81 = \qty{1.57}{s}$;
$v = 5.0 + 9.81 \times 1.57 = \qty{20.4}{m/s}$. Dengan \omterm{def:g10:energy-conservation:energy}{energi}:
$v = \sqrt{5.0^2 + 2 \times 9.81 \times 20} = \qty{20.4}{m/s}$ ---
sama saja.
\end{solution}

\begin{solution}{exo:g12:projectile-motion:8}
$t_f = 2 \times 40 \times \sin\ang{30}/9.81 = \qty{4.1}{s}$;
$h = (40\sin\ang{30})^2/(2 \times 9.81) = \qty{20.4}{m}$;
$R = 40^2 \sin\ang{60}/9.81 = \qty{1.4e2}{m}$.
\end{solution}

\begin{solution}{exo:g12:projectile-motion:9}
$R = 10.5^2\sin\ang{40}/9.81 \approx \qty{7.2}{m}$. Yang ditinggalkan: \omterm{def:g11:forces-and-motion:com}{pusat
massanya} lepas landas tinggi (bendanya terbentang) lalu mendarat rendah dan
jauh ke depan (kakinya dilempar ke muka), dan pelompatnya bukanlah titik ---
\qty{1.7}{m} yang hilang itu bersemayam pada ilmu ukurnya, bukan pada parabolanya.
\end{solution}

\begin{solution}{exo:g12:projectile-motion:10}
Bersamaan: keduanya memiliki $v_y(0) = 0$, sehingga waktunya sama
$t = \sqrt{2 \times 1.2/9.81} = \qty{0.49}{s}$; yang ditembakkan mendarat
$300 \times 0.49 \approx \qty{1.5e2}{m}$ jauhnya. Yang diuji adalah
kebebasan gerak mendatar dan gerak tegaknya.
\end{solution}

\begin{solution}{exo:g12:projectile-motion:11}
$\sin\bigl(2(\ang{90}-\alpha)\bigr) = \sin(\ang{180}-2\alpha) =
\sin(2\alpha)$: $R$ yang sama. Di sini $\sin(2\alpha) = gR/v_0^2 =
9.81 \times 50/625 = 0.785$: $\alpha = \ang{25.9}$ atau \ang{64.1}.
\end{solution}

\begin{solution}{exo:g12:projectile-motion:12}
$R_{\max} = v_0^2/g = \qty{6.5}{m}$ (pada \ang{45});
$h_{\max} = v_0^2/2g = \qty{3.3}{m}$ (lurus ke atas).
$\sin(2\alpha) = 9.81 \times 5.0/64 = 0.766$: $\alpha = \ang{25.0}$
atau \ang{65.0}.
\end{solution}

\begin{solution}{exo:g12:projectile-motion:13}
$y(15) = 15\tan\ang{35} - \dfrac{9.81 \times 15^2}
{2 \times 18^2 \cos^2\ang{35}} = 10.5 - 5.1 = \qty{5.4}{m}$:
bolanya melewati pagar \qty{4.0}{m} itu sekitar \qty{1.4}{m}.
\end{solution}

\begin{solution}{exo:g12:projectile-motion:14}
(a) \omterm{def:g11:circuits-and-power:resistance}{Hambatannya} mengimbangi \omterm{def:g10:universal-gravitation:weight}{beratnya}: $80 \times 9.81 \approx \qty{7.8e2}{N}$.
(b) $v(t)$ menanjak dengan kemiringan yang kian mengecil lalu mendatar pada
\qty{55}{m/s}. (c) Di ruang hampa:
$\sqrt{2 \times 9.81 \times 2000} \approx \qty{2.0e2}{m/s}$, dua puluh
kali \qty{9}{m/s} yang sesungguhnya --- tetes hujan menghabiskan hampir
seluruh jatuhnya pada \omterm{def:g12:projectile-motion:terminal}{laju terminal}.
\end{solution}

\begin{solution}{exo:g12:projectile-motion:15}
Pada $t^* = D/(v_0\cos\alpha)$ panahnya berada pada ketinggian
$v_0\sin\alpha\, t^* - \tfrac12 g t^{*2} = D\tan\alpha - \tfrac12 g
t^{*2} = H - \tfrac12 g t^{*2}$, persis $y = H - \tfrac12 g t^{*2}$
milik monyet yang jatuh itu. Keduanya menggantung $\tfrac12 gt^2$ di bawah
tempat tanpa-gravitasinya --- keduanya bertemu untuk $v_0$ berapa pun.
\end{solution}

\begin{solution}{pb:g12:projectile-motion:1}
\textbf{1.} Hanya \omterm{def:g10:universal-gravitation:weight}{beratnya} (udaranya diabaikan); \omterm{thm:g12:newtons-laws:second}{hukum kedua Newton}:
$m\vect a = m\vect g$, sehingga $\vect a = (0,\,-g)$.

\textbf{2.} $v_x = v_0\cos\alpha$, $v_y = v_0\sin\alpha - gt$.

\textbf{3.} $x = v_0\cos\alpha\, t$,
$y = v_0\sin\alpha\, t - \tfrac12 gt^2$.

\textbf{4.} $t = x/(v_0\cos\alpha)$:
$y = x\tan\alpha - g x^2/(2v_0^2\cos^2\alpha)$.

\textbf{5.} $v_0^2 = \dfrac{g d^2}{2\cos^2\alpha\,(d\tan\alpha -
1.05)} = \dfrac{173.0}{0.826 \times 3.96} = 52.9$:
$v_0 = \qty{7.28}{m/s} \approx \qty{26}{km/h}$.

\textbf{6.} $t = d/(v_0\cos\alpha) = 4.20/4.68 = \qty{0.90}{s}$.

\textbf{7.} Puncaknya pada $t = v_0\sin\alpha/g = \qty{0.57}{s}$, dengan tinggi
$2.00 + (v_0\sin\alpha)^2/2g = 2.00 + 1.58 = \qty{3.58}{m}$;
$0.57 < 0.90$: sudah menurun di ringnya.

\textbf{8.} $y(0.90) = 0.90\tan\ang{50} - 0.182 = \qty{0.89}{m}$
di atas titik lepasnya, yaitu \qty{2.89}{m} di atas lantainya: \qty{0.29}{m}
di atas ujung jarinya.

\textbf{9.} $v_x = \qty{4.68}{m/s}$;
$v_y = 5.57 - 9.81 \times 0.90 = \qty{-3.24}{m/s}$;
$v = \qty{5.69}{m/s}$, pada $\arctan(3.24/4.68) = \ang{34.7}$ di bawah
garis mendatarnya.

\textbf{10.} Ringnya berupa cincin mendatar: bolanya harus melintasi
bidangnya dari atas. Kian curam masuknya, kian besar bukaan cincinnya
tampak sepanjang kecepatannya --- kian lapang sisa untuk keliru.

\textbf{11.} $1.05/26.0 = 4\%$: terabaikan. $y(t_f) = 0$ memberikan
$t_f\,(v_0\sin\alpha - \tfrac12 g t_f) = 0$, sehingga
$t_f = 2v_0\sin\alpha/g$.

\textbf{12.} $R = v_0\cos\alpha\, t_f = v_0^2\sin(2\alpha)/g$;
terbesar pada $\alpha = \ang{45}$.

\textbf{13.} $v_0 = \sqrt{gR} = \sqrt{9.81 \times 26.0} =
\qty{16.0}{m/s} \approx \qty{57}{km/h}$.

\textbf{14.} $R_{\max} = 17.5^2/9.81 = \qty{31.2}{m} >
\qty{26.0}{m}$: terjangkau, dengan sisa \omterm{def:g12:projectile-motion:range}{jangkauan} \qty{5}{m}.

\textbf{15.} $\sin(2\alpha) = 255.1/306.3 = 0.833$:
$\alpha = \ang{28.2}$ atau \ang{61.8}; \omterm{def:g12:projectile-motion:range}{waktu terbangnya}
$2v_0\sin\alpha/g = \qty{1.69}{s}$ dan \qty{3.14}{s}. Hanya tembakan yang
datar itu yang mendahului bel \qty{2.0}{s}.

\textbf{16.} $v_0^2 = \dfrac{9.81 \times 36}{2\cos^2\ang{75}\,
(6.0\tan\ang{75} - 1.05)} = \dfrac{353.2}{2.86} = 123.5$:
$v_0 = \qty{11.1}{m/s}$.

\textbf{17.} $y(5.55) = 5.55\tan\ang{75} - 18.26 = \qty{2.45}{m}$
di atas titik lepasnya, yaitu \qty{4.45}{m} di atas lantainya: \qty{0.50}{m}
di atas tepi atas papannya.

\textbf{18.} Puncaknya $2.00 + (11.1\sin\ang{75})^2/2g \approx
\qty{7.9}{m}$; penerbangannya $6.0/(11.1\cos\ang{75}) = \qty{2.1}{s}$ ---
dua sekon penuh dengan bolanya menggantung \qty{8}{m} di atas.

\textbf{19.} Hambatan udara memendekkan setiap tembakan lalu mencuramkan
turunnya; pengaruhnya tumbuh bersama lajunya, sehingga lemparan sepanjang
lapangan \qty{17.5}{m/s} itu yang paling menderita --- \omterm{def:g12:projectile-motion:range}{jangkauan}
sesungguhnya kurang beberapa \omterm{def:g10:orders-of-magnitude:unit}{meter} dari \qty{31}{m}, memakan sisanya.

\textbf{20.} \emph{Putusan tembakan di detik terakhir}: tembakan
\qty{26.0}{m} itu menuntut \qty{16.0}{m/s} sementara sebuah lengan memasok
sekitar \qty{17.5}{m/s}, sisa \omterm{def:g12:projectile-motion:range}{jangkauan} \qty{5}{m} ($\approx 20\%$) yang
dipangkas hambatan udara tetapi tidak dihapusnya. Mungkin --- dan itulah
tepatnya sebabnya, beberapa kali tiap musim, bolanya masuk.
\end{solution}
